\documentclass[12pt,a4paper]{article}

\usepackage[a4paper,margin=1in]{geometry}
\usepackage[colorlinks=true,
linkcolor=KimPink,
urlcolor=KimPink]{hyperref}

\usepackage{setspace}
\usepackage{xcolor}
\usepackage{titlesec}
\usepackage{graphicx}
\usepackage[most]{tcolorbox}
\usepackage{tocloft}
\usepackage{titlesec}


\renewcommand{\cftsecfont}
{\color{DarkGrey}}

\renewcommand{\cftsecpagefont}
{\color{DarkGrey}}

\definecolor{KimPink}{HTML}{E75480}
\definecolor{DarkGrey}{HTML}{2D2D2D}

\titleformat{\section}
{\LARGE\bfseries\color{KimPink}}
{\thesection}
{1em}
{}
[\color{KimPink}\titlerule]

\titlespacing*{\section}
{0pt}
{24pt}
{12pt}
\titleformat{\subsection}
{\Large\bfseries\color{KimPink}}
{\thesubsection}
{1em}
{}
\titlespacing*{\subsection}
{0pt}
{14pt}
{8pt}
\pagestyle{plain}

\title{Computer Science Portfolio Report}
\author{Kimberly Njoki Gachuhi}
\date{June 2026}
\onehalfspacing

\setlength{\parskip}{1em}
\setlength{\parindent}{0pt}
\begin{document}


\begin{titlepage}
\thispagestyle{empty}

\begin{center}

\vspace*{2cm}

{\Huge\bfseries
Computer Science Portfolio Report
}

\vspace{0.3cm}

{\Large\color{KimPink}
Portfolio Website, GitHub Repository and Professional CV
}

\vspace{0.3cm}

\vspace{0.2cm}

\begin{tcolorbox}[
colback=KimPink,
colframe=KimPink,
height=0.8cm,
arc=0mm,
boxrule=0pt]
\end{tcolorbox}
\vspace{1.5cm}

{\Large Kimberly Njoki Gachuhi}

\vspace{0.5cm}

Bachelor of Science in Computer Science

\vspace{0.3cm}

GISMA University of Applied Sciences

\vspace{0.3cm}

June 2026
\vspace{1cm}

\begin{tcolorbox}[
colback=KimPink!5,
colframe=KimPink,
width=0.75\textwidth]

Building a professional portfolio has allowed me to showcase my skills, projects, and growth as an aspiring Full Stack Web Developer.

\end{tcolorbox}

\vfill


\vspace{1cm}



\end{center}

\end{titlepage}

\clearpage



\vspace{1cm}

\tableofcontents


\clearpage

\section{Introduction}
\begin{tcolorbox}[
colback=KimPink!5,
colframe=KimPink,
title=About This Report]

This report documents the development of my portfolio website, GitHub repository, CV, and technical projects completed during my Computer Science studies.

\end{tcolorbox}

My goal is to become a Full Stack Web Developer and build modern web applications that solve real problems while providing an excellent user experience. As a Computer Science student at GISMA University of Applied Sciences, I wanted to create a professional portfolio that showcases my technical skills, projects, academic background, and career development.

To guide the development of this portfolio, I selected a real Full Stack Web Developer job opportunity that reflects the type of role I aim to pursue in the future.






\textbf{Target Full Stack Web Developer Position with link provided below:}

link-[\href{https://www.linkedin.com/jobs/view/4419844638/?trackingId=PuSz2FWbfsMHEbID0iAecQ%3D%3D&refId=HBHbjzF3obmvOUfJKJ%2FX3A%3D%3D&eBP=NON_CHARGEABLE_CHANNEL&alternateChannel=search&isJobSearch=false&lipi=urn%3Ali%3Apage%3Ad_flagship3_nlsearch_srp_jobs%3B98OjX5xtQEyA8qJpKWoasw%3D%3D}{Senior Full Stack Developer} ]

The position requires skills in web development, programming, version control, problem solving, and collaboration, all of which influenced the content and design of my portfolio.The specifics are on React, Java Script, SQL, CSS and Docker or any equivalent.

The following resources were created as part of this project:

\begin{itemize}
\item GitHub Repository: \url{https://github.com/KimberlyNjoki/ComputerScienceProject}
\item Portfolio Website: \url{https://kimberlynjoki.github.io/PortfolioWebsite/}

\item Professional CV:
\href{https://raw.githubusercontent.com/KimberlyNjoki/ComputerScienceProject/main/LaTeXCV/LatexCV.pdf}{Download My CV}
\end{itemize}
\section{Portfolio, Repository and CV Structure}

To showcase my work and support my career goals, I created a portfolio website, a GitHub repository, and a professional CV. Each of these serves a different purpose but together they provide a complete overview of my skills, projects, and academic background.

\subsection{Portfolio Website Structure}

I designed the portfolio website as a single-page website because I wanted visitors to find everything in one place without having to move between different pages.

The website includes:

\begin{itemize}
\item \textbf{Hero Section} introducing who I am and what I am studying.
\item \textbf{About Section} sharing more about my background, interests, and career goals.
\item \textbf{Skills Section} highlighting the technologies and tools I have learned and used.
\item \textbf{Projects Section} showcasing projects I have worked on, including the Inizio Cosmetics platform, Todo Application, and Internship Tracker.
\item \textbf{Contact Section} providing ways to connect with me through email, GitHub, LinkedIn, and my downloadable CV situated at the top of the website.
\end{itemize}

\subsection{GitHub Repository Structure}

I used GitHub to keep all my project files organised and in one place.

The repository contains:

\begin{itemize}
\item \texttt{PortfolioWebsite/} containing the portfolio website files.
\item \texttt{LaTeXCV/} containing my CV and supporting files.
\item \texttt{LaTeXReport/} containing this project report.
\item \texttt{TechnicalExercises/} containing the E-Commerce Platform, Internship Tracker, and Todo Application projects.
\item \texttt{README.md} providing an overview of the repository.
\end{itemize}

GitHub also helped me keep track of changes as I continued developing and improving my work.

\subsection{CV Structure}

I created my CV using LaTeX to give it a clean and professional appearance.

The CV includes:

\begin{itemize}
\item Personal and contact information.
\item A short professional profile.
\item Technical skills.
\item Educational background.
\item Featured projects.
\item Languages spoken.
\item Career goals and interests.
\item Links to my portfolio website and GitHub repository.
\end{itemize}

Together, these three components give a clear picture of my work, skills, and progress as I continue working towards a career in full stack web development.


\section{Design Decisions and Justification}

The design of the portfolio website, CV, and project report was influenced by my goal of creating a professional profile that represents both my skills and my career aspirations as a Computer Science student.

For the portfolio website, I chose a dark theme with pink accent colours because I wanted the website to feel modern while still reflecting my personal style. I also wanted it to stand out from many portfolio websites that often use similar layouts and colour schemes. A single-page structure was selected to keep navigation simple and allow visitors to view all important information in one place. Project cards and organised sections were used to present information clearly and showcase my projects in a coherent format.

For the CV, I chose a clean LaTeX design because it creates a professional and well-structured document. The layout was organised to highlight my education, technical skills, projects, and career goals. I also included subtle pink accents in some headings to maintain consistency with my personal branding. Links to my portfolio website and GitHub repository were added to allow employers and recruiters to learn more about my work. A professional photograph was also included to make the CV more personal.

For the project report, I wanted the design to remain professional while matching the overall style of the portfolio website and CV. Consistent formatting and clear section headings were used to improve readability and help explain the development process more effectively.I also chose to add some pink accents to fit my personal brand.

Overall, these design choices helped me create a consistent and professional representation of my work while highlighting the skills and experience I have gained through my studies and personal projects.

\section{Technologies and Tools Used}

Different technologies and tools were used throughout the development of the portfolio website, CV, and project report.

For the portfolio website, I used HTML and CSS to build and style the website. Visual Studio Code was used to write and edit the code, while Git and GitHub were used to manage changes and publish the website online through GitHub Pages.

For the CV, I used LaTeX because I wanted a clean and professional document that would stand out from a standard word-processed CV. The CV was created and edited in Overleaf before being added to my GitHub repository.

For this project report, I used Microsoft Word to organise my ideas and plan the structure of the report before writing the final version in LaTeX. LaTeX was then used to produce a well-formatted academic document with consistent styling and layout.

Throughout the project, GitHub played an important role in storing files, tracking progress, and organising all project components in a single repository. Together, these tools helped me develop and present my work in a professional and organised manner.

\section{Strengths of My Portfolio}

One of the biggest strengths of this project is that it brings together everything I have worked on in one place. Instead of having my projects, CV, and report separated, they are all connected and easy to access through my portfolio and GitHub repository.

I am particularly happy with the portfolio website because it gives me a platform to showcase who I am, what I have learned, and the projects I have completed so far. The layout is simple, easy to navigate, and allows visitors to quickly find information about my skills, projects, and career interests.

Another strength is my CV. Creating it in LaTeX helped me produce a clean and professional document that clearly presents my education, skills, and projects. Including links to my portfolio website and GitHub repository also allows someone reading my CV to learn more about my work.

I also consider this report a strength because it documents the thought process behind the project and explains how everything was developed. It provides context for the portfolio and shows the work that took place behind the final result.

The technical projects included in my repository are another strong point. The Todo Application, Internship Tracker, and E-Commerce Platform helped me apply the skills I have learned and gave me practical projects to include in my portfolio.

Another strength is that the portfolio reflects many of the skills required for the Senior Full Stack Developer role I selected. My projects, GitHub repository, CV, and portfolio website demonstrate a strong foundation in web development, version control, problem solving, and project organisation.


Overall, I believe the strongest part of this project is that it reflects my progress as a Computer Science student and provides a professional way of presenting my skills, projects, and career goals.


\section{Weaknesses and Limitations}

Although I am proud of the final result, there are still areas I would like to improve.

One limitation is that I currently have a small number of projects in my portfolio. The projects included demonstrate my skills, but I would like to add more projects as I continue learning and gaining experience.

I also feel that my portfolio could benefit from more full stack projects. Most of my current work focuses on frontend development, and I would like to build more applications that include databases, authentication, and backend functionality.

The design can also continue to improve. While I am happy with the overall appearance, I know that as my web development skills grow, I will be able to create more polished and advanced user interfaces.

One skill I am currently developing is Docker. Although I have limited experience with containerisation at the moment, I am actively learning it and plan to use it in future projects to strengthen my full stack development skills.


For the CV, I would like to include more professional experience, certifications, and projects as I continue building my career. At the moment, it reflects my current stage as a student, but I expect it to grow over time.

Overall, these limitations reflect areas where I still want to grow and improve as I continue working towards a career in full stack web development.

\section{Future Improvements}

Despite being happy with what I have achieved so far, I see this project as the beginning rather than the final version.

One of my main goals is to add more projects to my portfolio as I continue learning and gaining experience. I want the portfolio to grow alongside my skills and better reflect the type of developer I am becoming.

I would also like to build more full stack applications with backend functionality, databases, and user authentication. This will help me strengthen my skills and prepare for the type of roles I hope to work in after university like the one mentioned in the beginning.

For the portfolio website, I would like to add a working contact form and make it more interactive. I am also considering adding a blog section where I can document projects, share what I am learning, and track my progress.

As I complete more projects, gain certifications, and build professional experience, I will continue updating my CV so that it reflects my latest achievements and skills.

My long-term goal is for this portfolio, CV, and repository to grow into a strong professional profile that helps me secure internships and eventually a career as a Full Stack Web Developer like the one I earlier stated.

\section{Conclusion}

Working on this project allowed me to bring together the skills I have developed throughout my studies and apply them to something meaningful. Creating the portfolio website, CV, GitHub repository, and project report helped me gain more confidence in my technical abilities and better understand how to present my work professionally.

The portfolio is an important step towards my goal of becoming a Full Stack Web Developer because it gives me a place to showcase my skills, projects, and progress. It also provides a professional way for potential employers and recruiters to learn more about me and the work I have done.

I believe this project reflects where I am currently in my learning journey. It highlights the skills I have gained so far while also showing the areas where I still want to grow and improve.

Although the project is complete, I do not see it as finished. I plan to continue adding new projects, improving my skills, and updating my portfolio and CV as I gain more experience. For me, this project is the beginning of building a professional profile that will grow alongside my career.

\begin{tcolorbox}[
colback=KimPink!5,
colframe=KimPink,
title=Final Reflection]

This portfolio project represents an important milestone in my journey as a Computer Science student and aspiring Full Stack Web Developer. It provides a strong foundation that I will continue improving throughout my studies and future career.

\end{tcolorbox}
\end{document}
